\section{Method}
\label{sec:method}

The design goal was that source code be the \emph{only} thing differing between
the binaries we compare. Everything below follows from that.

\subsection{Sandboxes and build}

Three clones of the GraphViz repository were made, all checked out to commit
\texttt{23b06521}. One was kept untouched as the baseline; the other two were
given to the agents. Each clone was placed on its own named branch so that the
agents' commit histories are preserved as artifacts rather than left on a
detached \texttt{HEAD}.

All three were built with identical CMake flags
(\texttt{-DCMAKE\_BUILD\_TYPE=Release}, Poppler and GVEdit disabled) and
installed to separate prefixes. Each resulting \texttt{dot} binary resolves its
libraries through an \texttt{@executable\_path}-relative rpath, so a binary
always loads its own build of the layout libraries. We verified this rather than
assuming it.

\subsection{The prompt}

Both agents received the same text. It stated the goal --- make \texttt{dot}
faster on large graphs without changing its output --- and deliberately named no
function, file, or phase. Finding the hot code was part of the task under study.
The prompt also:

\begin{itemize}
\item described the target workload shape (large layered graphs with heavy edge
      crossings) without supplying the evaluation graphs, so agents generated
      their own test inputs;
\item required byte-identical \texttt{-Tdot} output and pointed at the 265
      graphs in the repository's own \texttt{tests/} directory;
\item enumerated disallowed shortcuts: detecting or special-casing specific
      inputs, caching layout results across runs, reducing iteration counts or
      loosening convergence thresholds, hardcoding computed values, and
      modifying the test suite;
\item stated that measurement would be independent, that identical builds vary
      by several percent on the target machine, and that a modest honest result
      or no result at all was an acceptable outcome.
\end{itemize}

The prohibition on reducing iteration counts matters most. Crossing minimization
is iteration-bounded, so cutting its iteration limit produces a large, real, and
reproducible speedup that is not an optimization. The prompt reframed that as a
finding to report rather than a change to ship, giving an agent that noticed it
somewhere honest to put it.

Agents worked autonomously with no turn or time limit and were asked to commit
incrementally and to write a report describing what they profiled, what they
changed, what failed, and what they verified. Neither agent had access to the
evaluation harness, the benchmark graphs, or the other agent's sandbox.

\subsection{Measurement}

All timings use \texttt{hyperfine} 1.20.0~\cite{peter2023hyperfine} with 10
warmup and 10 timed runs.
Every binary for a given (graph, phase) is measured in a \emph{single}
\texttt{hyperfine} invocation, so the binaries run seconds apart under the same
thermal conditions. This matters: measuring them in separate invocations
produced order-dependent bias on our machine large enough to swamp the effects
we are trying to detect. Consequently, only comparisons made within one
invocation are meaningful, and we do not compare numbers across runs.

Phase attribution uses \texttt{-Gphase=N}, which runs a prefix of the pipeline
and returns. Phase timings are therefore \emph{cumulative}: phase~2 includes
phase~1. Per-stage costs are obtained by differencing adjacent phases, which
compounds the error of both terms; we report them as weaker evidence than the
cumulative figures throughout.

Evaluation graphs are layered directed graphs generated by a fixed-seed script:
100 nodes per rank, each node with three edges to randomly chosen nodes in the
next rank, at 5{,}000 and 10{,}000 nodes. The generator is deterministic --- we
verified that regenerating a graph reproduces it byte-for-byte --- so the
described test suite and the measured one cannot diverge.

Experiments ran on an Apple M4 Pro (12 cores: 8 performance, 4 efficiency),
24\,GB unified memory, macOS 26.5.2, with Apple Clang 17.

\subsection{The null test}
\label{sec:nulltest}

Before measuring any agent's work we established what ``no difference'' looks
like. The baseline build tree was installed to three separate prefixes,
producing three binaries that are byte-identical (SHA-1
\texttt{b02f718d\ldots}); these were then benchmarked against one another exactly
as the real binaries are, as three commands in one \texttt{hyperfine} invocation.

Any spread observed there is attributable to run ordering, thermal drift, and
per-run variance --- everything the real comparison measures except source code.
Section~\ref{sec:results} reports the result, and we use it as the floor below
which no claimed speedup is meaningful. We rebuilt nothing for this test: a
recompilation would embed a new build-date string into the version, introducing
a difference rather than holding source constant.

\subsection{Correctness checking}

Correctness is checked on \texttt{dot -Tdot} output rather than a rendered
image. The \texttt{dot} text format is deterministic, so byte-identity is a
meaningful assertion, whereas image comparison would fold in rendering and
floating-point behaviour we are not studying. We additionally read both agents'
diffs against the list of disallowed shortcuts in the prompt.
